% Appendix E -- Exchanging Documents with InitechDOS Programs (the 1991
% Macintosh User's Guide Appendix B "Exchanging Files and Disks With MS-DOS
% and Apple II Computers", apple-macintosh.md Sec 8.2, cut to the one disk
% format FLAIR uses).  Sources (Law 1): os/flair/stdfile.c and
% os/milton/kmain.c sf_name_ok (8.3 names, letters stored in capitals);
% os/apps/itext/itext.c (doc_load: CR LF and LF read as one line end;
% doc_write: CR line ends; FLAIR_TE_TEXT_MAX); os/milton/kmain.c
% finder_surface_open (documents are not opened from the desktop); the data
% volume (README.TXT has LF line ends: l_readme shows it read correctly).
% GATES: test-flair-itext test-itext-route test-stdfile test-type test-pubs-flair-screens
\Gates{test-flair-itext, test-itext-route, test-stdfile, test-type,
  test-pubs-flair-screens}
\chapter{Exchanging Documents with InitechDOS Programs}
\label{apx:exchange}
\index{InitechDOS programs!exchanging documents with}

FLAIR and the InitechDOS programs use the same disk and the same files.  A
document you save in FLAIR can be used at the InitechDOS prompt, and a text
file written by an InitechDOS program can be opened in FLAIR.  This
appendix tells you what to keep in mind.

\section{File names}
\index{file names}
Every file has an InitechDOS name: one to eight characters, then, if you
want, a period and one to three more.  FLAIR shows the name exactly as
InitechDOS stores it, in capital letters, and the Save dialog box turns the
letters you type into capitals.  A name that does not follow these rules
cannot be typed in the Save dialog box (Chapter~\ref{ch:appsref}).

Name your documents so that you can tell them apart at the InitechDOS
prompt as well as in FLAIR, for example with the extension \file{.TXT} for
text documents.

\section{Text documents}
\index{text files}\index{line ends}
InitechText saves plain text: the characters of the document, one byte
each, with a carriage return at the end of each line.  There are no type
styles or other formatting to lose.
\begin{itemize}
\item InitechDOS programs, and files written at the InitechDOS prompt,
  usually end each line with a carriage return and a line feed, or with a
  line feed alone.  InitechText reads either as one line end, so such a
  file opens correctly (Figure~\ref{fig:readme}).  If you save it from
  InitechText, its lines end with a carriage return alone.
\item InitechText opens text files of up to 4,096 characters.  A longer
  file gives the notice in Figure~\ref{fig:toolarge}.
\item The characters that can be typed and stored are listed in
  Appendix~C.
\end{itemize}

\section{Using a FLAIR document at the prompt}
\begin{itemize}
\item \file{TYPE} shows a text document on the screen
  (Figure~\ref{fig:ktype}).
\item \file{COPY} and \file{XCOPY} copy documents to another disk or
  directory; FLAIR itself moves documents but does not copy them.
\item A document saved in a FLAIR folder is in the directory of the same
  name: a document in the \file{APPS} folder is in \file{A:\textbackslash
  APPS}.
\end{itemize}

\section{Using a file from the prompt in FLAIR}
\begin{itemize}
\item A file you make at the prompt in a directory appears in the folder of
  the same name the next time its Finder window is opened.
\item Open a text file from within InitechText with \menuitem{Open...}.
  Double-clicking a document's icon does not open it.
\item Programs whose names end in \file{.COM} cannot be run from the
  desktop; run them at the prompt.
\end{itemize}

\section{Files to leave alone}
Do not change, rename or delete the file \file{DESKTOP.DB} or the
\file{TRASH} directory at the InitechDOS prompt, and do not move files in or
out of \file{TRASH} there (Chapter~\ref{ch:dos}).
